\documentclass[11pt]{article}
\usepackage[margin=1in]{geometry}
\usepackage{amsmath,amssymb,amsthm}
\usepackage{graphicx}
\usepackage[hidelinks]{hyperref}

\newtheorem{theorem}{Theorem}
\newtheorem{lemma}{Lemma}
\newtheorem{remark}{Remark}

\newcommand{\R}{\mathbb R}
\newcommand{\JF}{\mathrm{JF}}
\newcommand{\supp}{\mathrm{supp}}

\title{A Counterexample for the $\ell_\infty^d$-Ball Version of Problem 1}
\author{Run \texttt{fa\_banach\_001}, lane 0}
\date{2026-06-23}

\begin{document}
\maketitle

\begin{center}
\fbox{\parbox{0.9\linewidth}{
\textbf{Status.} Counterexample, likely valid.  For every nonempty bounded open
$\Omega\subset\R^d$, $d\ge2$, the James-function space defined using
$\ell_\infty^d$ balls contained in $\Omega$ contains an isomorphic copy of
$\ell_1$.  Thus the $\ell_\infty^d$-ball branch of Problem 1(i) in
arXiv:1210.2379 has a negative answer.  The Euclidean-ball
$\ell_1$-noncontainment branch is not settled here.
}}
\end{center}

\section{Source Problem}

The source paper is S. A. Argyros, A. Manoussakis, and M. Petrakis,
\emph{Function spaces not containing $\ell_1$}, arXiv:1210.2379 \cite{AMP}.
Problem 1 asks: for a family $Q(\Omega)$ of convex bodies contained in a
bounded open $\Omega\subset\R^{d_0}$, $d_0>1$, when does the space with norm
\[
\|f\|_{\JF(Q(\Omega))}
=
\sup\left\{\left(\sum_{i=1}^n\left(\int_{V_i} f\right)^2\right)^{1/2}:
\{V_i\}_{i=1}^n\subset Q(\Omega),\ V_i\cap V_j=\emptyset\right\}
\]
not contain $\ell_1$, and when is its dual nonseparable?  The authors note
that their parallelepiped proof does not yield similar results for other
families and state that the problem is open if $Q(\Omega)$ is either the
Euclidean balls or the $\ell_\infty^{d_0}$ balls contained in $\Omega$.

\begin{figure}[h]
\centering
\includegraphics[width=\linewidth]{figures/open_problem_crop.png}
\caption{Problem 1 and the ball-family remark in arXiv:1210.2379, PDF page 45.}
\end{figure}

\section{Proof Intuition}

The obstruction is that $\ell_\infty^d$ balls are axis-parallel cubes.  In
dimension at least two, a high-frequency one-dimensional derivative can be
spread across the remaining $d-1$ coordinates.  Integrating over a cube of side
length $s$ produces a factor $s^{d-1}$ from the passive coordinates and an
oscillatory primitive difference in the first coordinate.  This factor is
exactly what keeps the $\ell_2$ sum over disjoint cubes uniformly bounded.

At the same time, nested detector cubes inside the support cube read off
\[
t^{d-1}\sin(2\pi 2^n t).
\]
For every subsequence one can choose a single $t$ whose binary expansion makes
these sine values oscillate along a further subsequence.  Hence the bounded
sequence has no weak-Cauchy subsequence.  Rosenthal's $\ell_1$ theorem then
forces an $\ell_1$ subsequence.

\section{Counterexample}

Let $\Omega\subset\R^d$ be nonempty, bounded, and open, with $d\ge2$.  Let
$Q_\infty(\Omega)$ denote the family of open axis-parallel $\ell_\infty^d$
balls, equivalently open cubes with equal side lengths, contained in $\Omega$.

\begin{theorem}\label{thm:counterexample}
The space $\JF(Q_\infty(\Omega))$ contains an isomorphic copy of $\ell_1$.
Consequently, the $\ell_\infty^d$-ball case of Problem 1(i) in
\cite{AMP} has a negative answer.
\end{theorem}

\begin{proof}
Choose an open axis-parallel cube
\[
C=(a_1,a_1+L)\times\cdots\times(a_d,a_d+L)
\]
whose closure is contained in $\Omega$.  Set
\[
\lambda_n=2\pi 2^n
\]
and define $f_n\in L^1(\Omega)$ by
\[
f_n(x_1,\ldots,x_d)
=
\begin{cases}
\dfrac{\lambda_n}{L}
\cos\!\left(\lambda_n\dfrac{x_1-a_1}{L}\right),
& (x_1,\ldots,x_d)\in C,\\[6pt]
0, & \text{otherwise.}
\end{cases}
\]
We first show that $(f_n)$ is bounded in $\JF(Q_\infty(\Omega))$.

Let $V$ be an $\ell_\infty^d$ ball contained in $\Omega$, and let $s(V)$ be
its side length.  Since $V\cap C$ is a rectangle, writing $s=s(V)$ gives
\[
\left|\int_V f_n\right|
\leq
s^{d-1}\min\left\{2,\frac{\lambda_n}{L}s\right\}.
\]
Indeed, the product of the intersection lengths in the last $d-1$ coordinates
is at most $s^{d-1}$, and the first-coordinate primitive
\[
\sin\!\left(\lambda_n\frac{x_1-a_1}{L}\right)
\]
has oscillation at most $2$ and derivative bounded by $\lambda_n/L$.

Let $V_1,\ldots,V_m$ be pairwise disjoint members of $Q_\infty(\Omega)$, and
put $s_j=s(V_j)$.  Since the cubes are disjoint and contained in the bounded
set $\Omega$,
\[
\sum_{j=1}^m s_j^d\leq |\Omega|.
\]
Let $D<\infty$ be an upper bound for side lengths of cubes contained in
$\Omega$.  Split the indices according to $s_j\leq L/\lambda_n$ and
$s_j>L/\lambda_n$.

If $s_j\leq L/\lambda_n$, then, since $d\ge2$ and $\lambda_n\ge1$,
\[
\left|\int_{V_j}f_n\right|^2
\leq
\left(\frac{\lambda_n}{L}\right)^2s_j^{2d}
\leq
L^{d-2}\lambda_n^{2-d}s_j^d
\leq
M_1 s_j^d
\]
for a constant $M_1$ depending only on $L$ and $d$.  If $s_j>L/\lambda_n$,
then
\[
\left|\int_{V_j}f_n\right|^2
\leq
4s_j^{2d-2}
=4s_j^d s_j^{d-2}
\leq 4D^{d-2}s_j^d.
\]
Thus
\[
\sum_{j=1}^m\left|\int_{V_j}f_n\right|^2
\leq
\max\{M_1,4D^{d-2}\}\,|\Omega|,
\]
uniformly in $n$ and in the finite disjoint cube family.  Hence $(f_n)$ is
bounded in $\JF(Q_\infty(\Omega))$.

We next show that $(f_n)$ has no weak-Cauchy subsequence.  Let $(f_{n_k})$ be
an arbitrary subsequence.  Choose a further subsequence, still denoted
$(f_{m_j})$, such that $m_{j+1}-m_j\to\infty$ and $m_{j+1}-m_j\ge4$.

We claim that there is a $t\in(0,1)$ such that
\[
\sin(2\pi 2^{m_{2j-1}}t)\to 1,
\qquad
\sin(2\pi 2^{m_{2j}}t)\to -1.
\]
Choose the binary expansion of $t=0.\varepsilon_1\varepsilon_2\ldots$ as
follows.  For odd $j$, set
$(\varepsilon_{m_j+1},\varepsilon_{m_j+2})=(0,1)$; for even $j$, set
$(\varepsilon_{m_j+1},\varepsilon_{m_j+2})=(1,1)$.  Between these prescribed
blocks put zeros, except where a later prescribed block begins.  Since the
gaps $m_{j+1}-m_j$ tend to infinity, the fractional part of $2^{m_j}t$ tends
to $1/4$ along odd $j$ and to $3/4$ along even $j$, proving the claim.

For this fixed $t$, the cube
\[
C_t=(a_1,a_1+tL)\times\cdots\times(a_d,a_d+tL)
\]
belongs to $Q_\infty(\Omega)$.  Its integration functional is continuous on
$\JF(Q_\infty(\Omega))$, and
\[
\int_{C_t} f_n
=(tL)^{d-1}\sin(2\pi 2^n t).
\]
Along the subsubsequence $(m_j)$ these scalar values have two distinct
subsequential limits, $(tL)^{d-1}$ and $-(tL)^{d-1}$.  Therefore the original
subsequence $(f_{n_k})$ is not weakly Cauchy.  Since the subsequence was
arbitrary, $(f_n)$ has no weak-Cauchy subsequence.

By Rosenthal's $\ell_1$ theorem \cite{Rosenthal}, every bounded sequence in a
Banach space has either a weak-Cauchy subsequence or a subsequence equivalent
to the unit vector basis of $\ell_1$.  Applying this to the bounded sequence
$(f_n)$, we obtain a subsequence equivalent to the $\ell_1$ basis.  Hence
$\ell_1$ embeds into $\JF(Q_\infty(\Omega))$.
\end{proof}

\section{Verification and Limitations}

The proof is deterministic.  The boundedness estimate uses only the displayed
definition of the norm and the elementary side-length split at
$s=L/\lambda_n$.  The no-weak-Cauchy argument uses only cube integration
functionals and an explicit binary-expansion selection.  Rosenthal's theorem is
the standard final step converting a bounded sequence with no weak-Cauchy
subsequence into an $\ell_1$ basic subsequence.

This packet settles only the $\ell_\infty^d$-ball branch of the source remark,
and it settles that branch negatively for Problem 1(i).  It does not address
the Euclidean-ball $\ell_1$-noncontainment branch.  The earlier lane-0 partial
packet is preserved inside this packet under the \texttt{history/} subfolder,
in the directory named \texttt{1210.2379\_Q\_family\_nonseparable\_dual};
it proves
nonseparability of the dual for both Euclidean and $\ell_\infty^d$ ball
families.  Thus the current combined conclusion for $\ell_\infty^d$ balls is:
the dual is nonseparable, but the space contains $\ell_1$.

Before promotion, the run indexes were searched for arXiv:1210.2379, the
packet slug, and the phrases ``Euclidean balls,'' ``ell\_infty balls,'' and
``JF(Q).''  A bounded web search for arXiv:1210.2379 with ``Euclidean balls,''
``ell\_infty,'' and ``James function space'' found only the source paper and
no obvious later exact resolution of the stated ball-family problem.

\section{Human Review Recommendation}

The main review point is the interpretation of $\ell_\infty^d$ balls as
axis-parallel cubes.  Under that standard convention, the construction is
direct.  Reviewers should also check the binary-selection paragraph and the
uniform norm estimate, especially the use of $d\ge2$ in the small-cube side of
the split.

\begin{thebibliography}{9}
\bibitem{AMP}
S. A. Argyros, A. Manoussakis, and M. Petrakis,
\emph{Function spaces not containing $\ell_1$},
arXiv:1210.2379, 2012.

\bibitem{Rosenthal}
H. P. Rosenthal,
\emph{A characterization of Banach spaces containing $\ell_1$},
Proceedings of the National Academy of Sciences of the United States of America
71 (1974), 2411--2413.
\end{thebibliography}

\end{document}
